\section{Teacher and student}
\label{sec:teacher-and-student}

% Needs in the preamble of main.tex:
%   \usepackage{amssymb}
%   \usepackage{array}
%   \usepackage{tikz}
%   \usetikzlibrary{arrows.meta,positioning,calc}

This section introduces the teacher--student framework used in our paper, as
well as the diagnostic checks that follow from it. We use the addition task as
a running example throughout; Section~\ref{sec:setup} gives the details of the
tasks and the experimental setup.

\subsection{One round, two roles}
\label{sec:two-roles}

At the start of round $i$, we have a single model, $M_{i-1}$, alongside a pool of
hard problems (longer than anything the model has been trained on), with no
correct answers attached. We put $M_{i-1}$ to work in two distinct
roles (Figure~\ref{fig:roles}).

Acting as the \emph{student}, the model simply tackles a hard problem head-on
and produces a direct answer. Acting as the \emph{teacher}, it follows a
structured procedure: it breaks the hard problem into smaller, easier pieces,
solves each piece individually, stitches those partial solutions back together,
and decides whether to keep or discard the resulting answer in a final step we call
\emph{selection}. Crucially, the teacher is not a separate, stronger model. It
is the exact same model, just guided by a decomposition procedure.

The answers that survive selection form the training set $D_i$, which is then
used to update $M_{i-1}$ into $M_i$. Because the student is trained on the
teacher's outputs, each round simply pulls the student's predictions toward the
teacher's. Three practical consequences immediately follow:

\begin{itemize}
\item If the teacher produces the exact same answers as the student, there is
  nothing new to learn, and the round achieves nothing.
\item If the teacher's answers are wrong, training pulls the student toward
  errors, actively degrading performance.
\item If the teacher is fundamentally incapable of generating the solutions the
  student lacks, no amount of iteration will bridge the gap, and progress
  stalls.
\end{itemize}

For a round of self-improvement to actually work, the teacher's answers must
therefore \emph{differ} from the student's, be \emph{correct}, and
\emph{cover what the student is missing}. While the second and third requirements
restate the failure modes identified by \citet{shin2026compositional}, the first
is our own addition.

\begin{figure}[H]
\centering
\begin{tikzpicture}[
  font=\small,
  node distance=5mm and 5mm,
  box/.style={draw, rounded corners=2pt, minimum height=8mm,
    inner xsep=4pt, align=center},
  model/.style={box, fill=black!10},
  plain/.style={align=center},
  flow/.style={-{Stealth[length=2mm]}, semithick},
  chk/.style={font=\footnotesize\itshape, text=blue!55!black, align=center},
  chkline/.style={dashed, blue!55!black},
]
% Teacher row
\node[box] (split) {Split};
\node[model, right=of split] (pieces) {Answer\\each piece};
\node[box, right=of pieces] (combine) {Combine};
\node[box, right=of combine] (select) {Select};
\node[plain, right=of select] (built) {kept\\answer};
\node[box, right=7mm of built] (train) {Train};
\node[plain, right=of train] (next) {$M_i$};
% Student row
\node[model, below=13mm of pieces] (direct) {Answer\\directly};
\node[plain] (dans) at (built |- direct) {direct\\answer};
% Problem
\node[plain, left=6mm of split] (x) at ($(split.west)!0.5!(split.west |-
  direct)$) {hard\\problem};
% Flows
\draw[flow] (x.east) -- ++(2mm,0) |- (split.west);
\draw[flow] (x.east) -- ++(2mm,0) |- (direct.west);
\draw[flow] (split) -- (pieces);
\draw[flow] (pieces) -- (combine);
\draw[flow] (combine) -- (select);
\draw[flow] (select) -- (built);
\draw[flow] (built) -- (train);
\draw[flow] (train) -- (next);
\draw[flow] (direct) -- (dans);
% Role labels
\node[plain, below=1mm of split.south west, anchor=north west]
  {\textbf{Teacher}};
\node[plain, below=1mm of direct.south west, anchor=north west]
  {\textbf{Student}};
% Checks
\draw[chkline, {Stealth[length=1.6mm]}-{Stealth[length=1.6mm]}]
  (built) -- node[chk, right=1pt] {novelty:\\do they\\differ?} (dans);
\node[chk, above=4mm of combine] (agr)
  {agreement: does a second split give the same answer?};
\draw[chkline] (agr.west) -| (split.north);
\node[chk] (con) at ($(combine.south)+(0,-7.5mm)$)
  {construction: is the\\answer ever new?};
\draw[chkline] (con.north) -- (combine.south);
\end{tikzpicture}
\caption{One round of self-improvement, seen as two roles of one model. Shaded
  boxes are calls to the model $M_{i-1}$. The teacher's kept answers are the
  training set $D_i$. Dashed lines mark where each check looks; none of them
  uses a correct answer to the hard problem.}
\label{fig:roles}
\end{figure}

\subsection{A worked example}
\label{sec:worked-example}

Figure~\ref{fig:example} walks through a concrete example of both roles in
action. Here, a model trained only on three- to seven-digit numbers is asked to
add two eight-digit numbers. Unsurprisingly, the student's direct attempt fails.
To solve it as a teacher, the model splits each number in two, adds the left
and right pieces separately, and concatenates the two results. With the first
split, this breaks down: the right-hand pieces sum to a five-digit number, and
the extra carry digit gets stuck awkwardly in the middle of the combined string
instead of carrying over to the left. However, the next two splits happen to
cut at boundaries with no carry, so both yield the right total.

What enters the training set next comes down to the selection rule. Without any
filtering, the teacher blindly keeps the flawed first answer. The filter
proposed by \citet{shin2026compositional}, which we call the \emph{no-carry
rule}, recognizes that an oversized partial sum signals an overflow error and
discards the example entirely. By contrast, our \emph{agreement rule} requires
no domain-specific arithmetic knowledge: it simply keeps an answer if at least
two different splits arrive at the same result.

\begin{figure}[t]
\centering
\small
\begin{tabular}{@{}lll@{}}
\toprule
\multicolumn{3}{@{}l}{\textbf{Problem:} \texttt{35128624 + 21415273}
  \quad (correct answer \texttt{56543897}, not available to the model)} \\
\midrule
\multicolumn{3}{@{}l}{\textbf{Student} answers directly:
  \texttt{56533897} \;$\times$} \\
\midrule
\multicolumn{3}{@{}l}{\textbf{Teacher} splits, answers each piece, and writes
  the sums side by side:} \\
Split & Sums of the pieces & Built answer \\
\texttt{3512|8624} & \texttt{3512+2141 = 5653},\;
  \texttt{8624+5273 = 13897} & \texttt{565313897} \;$\times$ \\
\texttt{351|28624} & \texttt{351+214 = 565},\;
  \texttt{28624+15273 = 43897} & \texttt{56543897} \;\checkmark \\
\texttt{35128|624} & \texttt{35128+21415 = 56543},\;
  \texttt{624+273 = 897} & \texttt{56543897} \;\checkmark \\
\midrule
\multicolumn{3}{@{}l}{\textbf{Selection} decides what is kept for training:} \\
No selection & keeps the first split's answer & \texttt{565313897}
  \;$\times$ \\
No-carry rule  & \texttt{13897} does not fit in four digits & nothing kept \\
Agreement rule & two of the three splits agree & \texttt{56543897}
  \;\checkmark \\
\bottomrule
\end{tabular}
\caption{One addition problem seen by the student and by the teacher. The
  marks \checkmark{} and $\times$ are for the reader: the model never sees
  them. The student's answer is illustrative. The no-carry rule is the
  hand-written filter of \citet{shin2026compositional}; the agreement rule
  needs no knowledge of carrying.}
\label{fig:example}
\end{figure}

\subsection{Two questions before every round}
\label{sec:two-questions}

Before committing to a training round, we draw a fresh batch of hard problems
and have the model play both roles. No weights are updated, and no correct
answers are consulted. Instead, we ask two diagnostic questions.

\paragraph{Is there anything new?} The \emph{novelty check} calculates the
fraction of the teacher's kept answers that actually differ from the student's
direct predictions. If nearly all of them match what the student already outputs,
training will be a dead end. With a reliable teacher, novelty is high on problem
lengths the student has yet to master, and drops near zero on lengths it
already knows. While comparing direct and decomposed answers has previously
been used to catch reasoning errors \citep{bottman2026operadic,yang2024decc}, we
use it here to gauge whether a training round has anything left to teach. Of
course, novelty alone cannot distinguish progress from regression: a different
answer might fix a student's mistake, or introduce a brand-new one.

\paragraph{Does a second split agree?} A valid solution should not depend on
where we happened to slice the input. Combining errors, on the other hand, are
usually sensitive to the cut point, as illustrated in Figure~\ref{fig:example} where a
carry ruins the first split while leaving the others untouched. To test for
this, the \emph{agreement check} re-derives each kept answer using alternative
splits and measures the fraction reproduced identically by at least one other
split, which we call the \emph{confirmed share}. A low confirmed share warns
that much of the teacher's training data is unreliable. Two caveats are worth
noting here. First, the alternative splits must cut at distinct positions; two
splits sharing a boundary will fail at that boundary in the exact same way,
giving a false sense of consensus. Second, agreement does not guarantee
correctness: two different splits can still converge on the same mistake.

The agreement rule of Figure~\ref{fig:example} is this same check, used as a
selection step. Because it keeps only answers that a second split confirms,
any teacher using it has a confirmed share of 100\% by construction, and only
the novelty check remains informative in that setting.

\subsection{One question per procedure}
\label{sec:one-question}

Our third diagnostic evaluates the decomposition procedure itself, not any
particular round, so it only needs to be run once per procedure.

\paragraph{Can the teacher build anything new?} The \emph{construction check}
measures how often the teacher's final answer is identical to one of the
individual sub-problem answers it started with. If this happens 100\% of the
time, the teacher is merely \emph{handing back} partial solutions: every answer
it produces for a hard problem is already something the model outputs for an
easier piece, so it can show the student nothing beyond the sub-problems.
Concatenating sums never falls into this trap, because the combined number is
always longer than either partial sum. On our second task, finding the longest
run of repeated characters in a string, the procedure from
\citet{shin2026compositional} always hands back: it simply returns the longest
run found inside an isolated piece, meaning it can never detect a run spanning
across a cut. Because this limitation is structural, we can measure it upfront
by substituting an oracle solver for the model before any training begins. In
practice, real models yielded identical scores on addition and nearly identical
ones on the string task (Table~\ref{tab:check3}).

That said, passing the construction check is necessary rather than sufficient.
A procedure might be capable of constructing novel answers in general while
still systematically filtering out the exact cases the student needs to see.
The no-carry rule illustrates this blind spot: it constructs valid multi-digit
sums beyond the student's reach, yet because it discards every problem
involving a carry across the boundary, the student is never shown a carry at
that position. The construction check is blind to this filtering bias, a
limitation whose consequences we examine in Section~\ref{sec:three-further-findings}.

\subsection{From answers to a prediction}
\label{sec:prediction}

For any given round and problem length, we combine our two per-round checks to
forecast one of three outcomes: the round will \emph{pay off}, \emph{do
nothing}, or \emph{do harm}. The rule also includes a third per-round check that we specified in advance:
because the teacher relies on the model's ability to solve smaller
sub-problems, a \emph{basics check} measures accuracy on easy problems.
The prediction is given by three questions asked in order
(Figure~\ref{fig:prediction}); the first that is answered ``no'' decides.

% Flowchart of the prediction rule. Needs in the preamble:
%   \usepackage{tikz}
%   \usetikzlibrary{arrows.meta,positioning}
\begin{figure}[t]
    \centering
    \begin{tikzpicture}[
      font=\small,
      every node/.append style={execute at begin node={\hyphenpenalty=10000\relax}},
      node distance=5mm and 16mm,
      q/.style={draw, rounded corners=2pt, text width=70mm, align=left,
        inner sep=5pt},
      res/.style={draw, thick, rounded corners=2pt, text width=38mm,
        align=left, inner sep=5pt, fill=black!8},
      flow/.style={-{Stealth[length=2mm]}, semithick},
      lab/.style={font=\footnotesize\itshape, inner sep=2pt},
    ]
    \node[q] (basics) {\textbf{Basics.} Are the easy problems solved?\\
      {\footnotesize Accuracy is at least 98\% on every group of easy problems.}};
    \node[q, below=of basics] (novelty) {\textbf{Novelty.} Is there anything
      new?\\
      {\footnotesize The teacher keeps at least 5\% of the sample, and at least
      50\% of the kept answers differ from the student's.}};
    \node[q, below=of novelty] (agree) {\textbf{Agreement.} Does a second split
      agree?\\
      {\footnotesize At least 38\% of the kept answers are confirmed by another
      split.}};
    \node[res, right=of basics] (harm1) {\textbf{Does harm}\\
      {\footnotesize The teacher would build on wrong answers.}};
    \node[res, right=of novelty] (nothing) {\textbf{Does nothing}\\
      {\footnotesize The teacher has nothing to teach.}};
    \node[res, right=of agree] (harm2) {\textbf{Does harm}\\
      {\footnotesize The teacher's answers cannot be trusted.}};
    \node[res, below=of agree.south west, anchor=north west] (pays)
      {\textbf{Pays off}\\
      {\footnotesize Enough new answers, and enough of them confirmed.}};
    \draw[flow] (basics) -- node[lab, above] {no} (harm1);
    \draw[flow] (novelty) -- node[lab, above] {no} (nothing);
    \draw[flow] (agree) -- node[lab, above] {no} (harm2);
    \draw[flow] ([xshift=10mm]basics.south west) --
      node[lab, right] {yes} ([xshift=10mm]novelty.north west);
    \draw[flow] ([xshift=10mm]novelty.south west) --
      node[lab, right] {yes} ([xshift=10mm]agree.north west);
    \draw[flow] ([xshift=10mm]agree.south west) --
      node[lab, right] {yes} ([xshift=10mm]pays.north west);
    \end{tikzpicture}
    \caption{How the checks give a prediction for one round at one problem length.
      The questions are asked in order, and the first ``no'' decides. Thresholds
      were calibrated on a separate set of pilot runs prior to evaluation (Section~\ref{sec:setup}).
      The basics check is the only one that uses correct answers, and only for easy problems.}
    \label{fig:prediction}
    \end{figure}

Novelty is tested before agreement because an answer the student already knows
can never change the model. Both the ordering of these rules and their numerical 
thresholds were calibrated on a separate set of pilot runs prior to evaluation (Section~\ref{sec:setup}).
The construction check is omitted from these round-by-round rules since it is
run once upfront to verify whether a procedure can ever teach anything beyond
the pieces.

Table~\ref{tab:checks} summarizes all four checks, and
Section~\ref{sec:do-the-checks-predict} evaluates how well our forecasts matched
actual training outcomes. In short, the novelty and agreement checks proved
highly predictive, whereas the basics check did not.

\begin{table}[t]
\centering
\caption{The checks. Kept answers are those the teacher's selection step
  retains for training. The first three use no correct answers. The version of
  the construction check for addition was defined after the evaluation runs.}
\label{tab:checks}
\small
\begin{tabular}{@{}l>{\raggedright\arraybackslash}p{0.25\linewidth}>{\raggedright\arraybackslash}p{0.36\linewidth}l@{}}
\toprule
Check & Question & What is measured & Asked \\
\midrule
Novelty & Is there anything new? & Share of kept answers that differ from the
  student's direct answer & Every round \\
\addlinespace
Agreement & Does a second split agree? & Share of kept answers that another
  split reproduces & Every round \\
\addlinespace
Construction & Can the teacher build anything new? & Share of the teacher's
  answers identical to an answer for one of the pieces & Once per procedure \\
\addlinespace
Basics & Are the easy problems solved? & Accuracy on easy problems & Every
  round \\
\bottomrule
\end{tabular}
\end{table}